\section{v1.4.0: two-level reports and endpoint trends}
\subsection{One series, two coordinated panels}
\key{\gradsparkcompare} derives a left dumbbell from the first and last positions,
and preserves the complete series on the right. Both panels share one explicit
range and unit. A fixed range or named scale is required; only line type is supported.
\begin{demo}{Project scores: endpoints do not skip NA}
\begingroup
\gradscale{project-score}{min=0,max=100}
\gradbarssetup{width=65mm,height=5ex,
 precision=0,label width=20mm,compare label=delta}
\begin{tabular}{lc}
Project & First / last and complete trend\\
\hline
Atlas & \gradsparkcompare[scale group=project-score]
 {20,40,NA,80}\\
Birch & \gradsparkcompare[scale group=project-score]
 {NA,30,60,90}\\
Elm & \gradsparkcompare[scale group=project-score]
 {20,50,65,NA}\\
\end{tabular}
\endgroup
\end{demo}
NA retains its position. A missing endpoint omits that marker and the connector;
the other endpoint remains, but the delta is missing. Single-point input uses
the same observation twice; empty input is an error. Default \key{compare label=values}
prints first/last; \key{delta} prints last minus first in the same unit.
Explicit text overrides the label; percent comparison labels are rejected.

\key{width} includes both panels, the default 3mm gap and two outer allowances
of marker size plus 1pt. \key{compare panel ratio=.4} divides the remaining width;
it must be strictly between zero and one, and the gap must be nonnegative.
Each panel must exceed twice the allowance. Outside labels reserve extra space.
Each composition produces one Figure with full series and endpoint alternatives;
explicit alt, artifact/off and embedded TikZ retain existing rules.
\clearpage
\subsection{Two explicit CSV text levels}
\key{group levels} accepts exactly two distinct declared text column IDs.
It is separate from old \key{group by} and manual groups; they cannot be mixed.
The illustrative CSV interleaves North and South; both have a Review team.
\begin{tcblisting}{enhanced,breakable,listing engine=listings,
 listing options={style=gradmanual},text and listing,
 title={Hierarchical project review, visible means and parent ranks},
 colback=gradbarsTrack!50,colbacklower=blue!3,colframe=gradbarsInk!25,
 coltitle=gradbarsInk,colbacktitle=gradbarsTrack,fonttitle=\bfseries,
 boxrule=.4pt,arc=0pt,left=3mm,right=3mm,top=2mm,bottom=2mm,
 before upper={\setlength{\parskip}{6pt}\small}}
\gradbarsloadcsv{projects}{docs/data/hierarchical-teams.csv}
\begin{center}
\begin{gradtable}[data=projects,style=hierarchical,
 width=\linewidth,group levels={region,team},
 rank scope=parent,parent summary=true,leaf summary=true,
 parent summary label={Region mean},
 leaf summary label={Team mean},label width=17mm]
\gradcolumn[source=Project,width=22mm]{project}{Project}
\gradcolumn[source=Region,width=17mm]{region}{Region}
\gradcolumn[source=Team,width=18mm]{team}{Team}
\gradcolumn[source=Score,type=number,width=20mm,
 aggregate=mean,mark=both,options={precision=1}]{score}{Score}
\gradcolumn[source=Trend,type=trendcompare,
 options={min=0,max=100,height=5mm,precision=0,
 compare label=delta}]{trend}{Endpoints / trend}
\gradsummary{All visible}
\end{gradtable}
\end{center}
\end{tcblisting}
Groups follow first appearance in the original CSV. Parent identity plus literal
leaf text forms a structured path, so same-named children remain separate.
Outer spaces are trimmed; empty keys use the missing heading label, whereas
literal NA is a distinct text key. Commas and slashes do not split paths.
\subsection{Selection, scope and native compound columns}
Selection remains filter, stable global sort, global limit, then grouping. Rows
retain the selected order within each leaf; groups retain source appearance order.
No implicit group limit or group scale is added.

\key{rank scope=table} remains default; \key{parent} ranks selected rows across
one parent's leaves, and \key{leaf} ranks within one path. Ties and the second
distinct value use canonical raw numbers. Old \key{rank scope=group} is rejected
for the new hierarchy. Summary rows never participate in ranks or scales.

\key{parent summary} and \key{leaf summary} default to false and can be enabled
independently. Each level has its own label and \key{scope=visible|data} via
\key{parent summary scope} or \key{leaf summary scope}; defaults are visible.
Data scope uses all bound rows in that node, even filtered or limited ones;
completely invisible groups still have no headings or summaries.
Column \key{aggregate} remains explicit. Overall \key{\gradsummary} has its own
scope. Legacy group summary keys are not used with two levels.

Native \key{type=trendcompare} accepts canonical braced lists or quoted CSV
series, defaults to 4ex height and resolves one range for both panels.
\key{scale=column} uses every valid observation, respecting table scale scope.
No scalar sort, filter, mark, aggregate or regional input applies to this compound
column. No arbitrary custom column adapter or tabularray helper is added.

\subsection{Long tables and layout limits}
\key{style=hierarchical} supplies accessible colors, row spacing and indented leaf
headings. In native \key{long=true} tables, parent/child headings stay with the
first row, and each summary chain with its last row; data rows can span pages.
An indivisible row/keep-chain must fit within a page. Booktabs, grid, rows and
none rule modes are checked against actual PDF shipout pages. This is no new
complete hierarchy Tagged PDF contract or PDF/UA certification.

Compile \path{docs/examples/hierarchical-longtable.tex} twice for the complete
72-row multipage report. Matching English and Chinese short examples are in
\path{docs/examples/}. Detailed contracts: \path{docs/hierarchical-tables.md}
and \path{docs/trend-compositions.md}.

Custom scalar backends require explicit TikZ fallback for \key{spark-compare};
backend schema remains 1.
